\documentclass[12pt]{article}
\usepackage{latexblog}

% ============================================================
% Blog Metadata
% ============================================================
\blogtitle{理解Java并发(3)：CAS}
\blogdate{2020-05-20}
\blogtags{Hotspot, JVM, 并发编程, Concurrent, ~理解Java并发, 刨根问底}
\bloglang{zh}
\blogsummary{}

\begin{document}
\makeblogtitle

CAS：Compare and Swap，即比较再交换。jdk5增加了并发包java.util.concurrent.*,其下面的类使用CAS算法实现了区别于synchronouse同步锁的一种乐观锁。

\section{实现机制}\label{ux5b9eux73b0ux673aux5236}

在Java中很多原子操作是通过CAS实现的，是一种无锁的方式，比如AtomicIntger的自增操作：

\begin{Shaded}
\begin{Highlighting}[]
\KeywordTok{public} \KeywordTok{class} \BuiltInTok{AtomicInteger} \KeywordTok{extends} \BuiltInTok{Number} \KeywordTok{implements}\NormalTok{ java}\OperatorTok{.}\FunctionTok{io}\OperatorTok{.}\FunctionTok{Serializable} \OperatorTok{\{}
    \KeywordTok{private} \DataTypeTok{static} \DataTypeTok{final} \DataTypeTok{long}\NormalTok{ serialVersionUID }\OperatorTok{=} \DecValTok{6214790243416807050L}\OperatorTok{;}

    \CommentTok{// setup to use Unsafe.compareAndSwapInt for updates}
    \KeywordTok{private} \DataTypeTok{static} \DataTypeTok{final}\NormalTok{ Unsafe unsafe }\OperatorTok{=}\NormalTok{ Unsafe}\OperatorTok{.}\FunctionTok{getUnsafe}\OperatorTok{();}
    \KeywordTok{private} \DataTypeTok{static} \DataTypeTok{final} \DataTypeTok{long}\NormalTok{ valueOffset}\OperatorTok{;}

    \DataTypeTok{static} \OperatorTok{\{}
        \ControlFlowTok{try} \OperatorTok{\{}
\NormalTok{            valueOffset }\OperatorTok{=}\NormalTok{ unsafe}\OperatorTok{.}\FunctionTok{objectFieldOffset}
                \OperatorTok{(}\BuiltInTok{AtomicInteger}\OperatorTok{.}\FunctionTok{class}\OperatorTok{.}\FunctionTok{getDeclaredField}\OperatorTok{(}\StringTok{"value"}\OperatorTok{));}
        \OperatorTok{\}} \ControlFlowTok{catch} \OperatorTok{(}\BuiltInTok{Exception}\NormalTok{ ex}\OperatorTok{)} \OperatorTok{\{} \ControlFlowTok{throw} \KeywordTok{new} \BuiltInTok{Error}\OperatorTok{(}\NormalTok{ex}\OperatorTok{);} \OperatorTok{\}}
    \OperatorTok{\}}

    \KeywordTok{private} \KeywordTok{volatile} \DataTypeTok{int}\NormalTok{ value}\OperatorTok{;}
    \CommentTok{// ...}
    \KeywordTok{public} \DataTypeTok{final} \DataTypeTok{int} \FunctionTok{getAndIncrement}\OperatorTok{()} \OperatorTok{\{}
        \ControlFlowTok{return}\NormalTok{ unsafe}\OperatorTok{.}\FunctionTok{getAndAddInt}\OperatorTok{(}\KeywordTok{this}\OperatorTok{,}\NormalTok{ valueOffset}\OperatorTok{,} \DecValTok{1}\OperatorTok{);}
    \OperatorTok{\}}
    \CommentTok{// ...}
\OperatorTok{\}}
\end{Highlighting}
\end{Shaded}

这里的操作调用了Unsafe类来实现，而这个Unsafe类里面的实代码是这样的：

\begin{Shaded}
\begin{Highlighting}[]
\KeywordTok{public} \DataTypeTok{final} \KeywordTok{native} \DataTypeTok{boolean} \FunctionTok{compareAndSwapInt}\OperatorTok{(}\BuiltInTok{Object}\NormalTok{ var1}\OperatorTok{,} \DataTypeTok{long}\NormalTok{ var2}\OperatorTok{,} \DataTypeTok{int}\NormalTok{ var4}\OperatorTok{,} \DataTypeTok{int}\NormalTok{ var5}\OperatorTok{);}
\KeywordTok{public} \DataTypeTok{final} \DataTypeTok{int} \FunctionTok{getAndAddInt}\OperatorTok{(}\BuiltInTok{Object}\NormalTok{ var1}\OperatorTok{,} \DataTypeTok{long}\NormalTok{ var2}\OperatorTok{,} \DataTypeTok{int}\NormalTok{ var4}\OperatorTok{)} \OperatorTok{\{}
    \DataTypeTok{int}\NormalTok{ var5}\OperatorTok{;}
    \ControlFlowTok{do} \OperatorTok{\{}
\NormalTok{        var5 }\OperatorTok{=} \KeywordTok{this}\OperatorTok{.}\FunctionTok{getIntVolatile}\OperatorTok{(}\NormalTok{var1}\OperatorTok{,}\NormalTok{ var2}\OperatorTok{);}
    \OperatorTok{\}} \ControlFlowTok{while}\OperatorTok{(!}\KeywordTok{this}\OperatorTok{.}\FunctionTok{compareAndSwapInt}\OperatorTok{(}\NormalTok{var1}\OperatorTok{,}\NormalTok{ var2}\OperatorTok{,}\NormalTok{ var5}\OperatorTok{,}\NormalTok{ var5 }\OperatorTok{+}\NormalTok{ var4}\OperatorTok{));}
    \ControlFlowTok{return}\NormalTok{ var5}\OperatorTok{;}
\OperatorTok{\}}
\end{Highlighting}
\end{Shaded}

而最终是用到了native的方法compareAndSwapInt，这部分的源码可以在OpenJDK源码中找到，而最终是会调用到\href{https://www.felixcloutier.com/x86/cmpxchg}{CMPXCHG}这样到CPU指令来完成CAS。

\begin{Shaded}
\begin{Highlighting}[]
\NormalTok{UNSAFE\_ENTRY}\OperatorTok{(}\NormalTok{jboolean}\OperatorTok{,}\NormalTok{ Unsafe\_CompareAndSwapInt}\OperatorTok{(}\NormalTok{JNIEnv }\OperatorTok{*}\NormalTok{env}\OperatorTok{,}\NormalTok{ jobject unsafe}\OperatorTok{,}\NormalTok{ jobject obj}\OperatorTok{,}\NormalTok{ jlong offset}\OperatorTok{,}\NormalTok{ jint e}\OperatorTok{,}\NormalTok{ jint x}\OperatorTok{))}
\NormalTok{  UnsafeWrapper}\OperatorTok{(}\StringTok{"Unsafe\_CompareAndSwapInt"}\OperatorTok{);}
\NormalTok{  oop p }\OperatorTok{=}\NormalTok{ JNIHandles}\OperatorTok{::}\NormalTok{resolve}\OperatorTok{(}\NormalTok{obj}\OperatorTok{);}
\NormalTok{  jint}\OperatorTok{*}\NormalTok{ addr }\OperatorTok{=} \OperatorTok{(}\NormalTok{jint }\OperatorTok{*)}\NormalTok{ index\_oop\_from\_field\_offset\_long}\OperatorTok{(}\NormalTok{p}\OperatorTok{,}\NormalTok{ offset}\OperatorTok{);}
  \ControlFlowTok{return} \OperatorTok{(}\NormalTok{jint}\OperatorTok{)(}\NormalTok{Atomic}\OperatorTok{::}\NormalTok{cmpxchg}\OperatorTok{(}\NormalTok{x}\OperatorTok{,}\NormalTok{ addr}\OperatorTok{,}\NormalTok{ e}\OperatorTok{))} \OperatorTok{==}\NormalTok{ e}\OperatorTok{;}
\NormalTok{UNSAFE\_END}
\end{Highlighting}
\end{Shaded}

其中，源码位于：

\begin{itemize}
\tightlist
\item
  \href{http://hg.openjdk.java.net/jdk8/jdk8/hotspot/file/tip/src/share/vm/prims/unsafe.cpp}{openjdk8/unsafe.cpp}
\item
  \href{http://hg.openjdk.java.net/jdk8/jdk8/hotspot/file/tip/src/share/vm/runtime/atomic.cpp}{openjdk8/atomic.cpp}
\end{itemize}

\section{ABA问题}\label{abaux95eeux9898}

使用CAS一个问题就是可能出现A-B-A问题，因为CAS到逻辑就是，假设更新到时候判断出值跟期望的一致那么就会进行更改；但实际值跟期望的一致并不能代表值没有发生过变化。可能的场景就是，值被改成别的然后又改回来了，就是从A-B，然后又变成A。

解决这个问题的思路在于，想办法标记出变化。通常的做法是利用版本号（而不是值），保证版本号每次都会变化。JDK中提供了AtomicStampedReference来解决A-B-A问题，实现如下：

\begin{Shaded}
\begin{Highlighting}[]
 \KeywordTok{public} \DataTypeTok{boolean} \FunctionTok{compareAndSet}\OperatorTok{(}\NormalTok{V   expectedReference}\OperatorTok{,}
\NormalTok{                              V   newReference}\OperatorTok{,}
                              \DataTypeTok{int}\NormalTok{ expectedStamp}\OperatorTok{,}
                              \DataTypeTok{int}\NormalTok{ newStamp}\OperatorTok{)} \OperatorTok{\{}
\NormalTok{     Pair}\OperatorTok{\textless{}}\NormalTok{V}\OperatorTok{\textgreater{}}\NormalTok{ current }\OperatorTok{=}\NormalTok{ pair}\OperatorTok{;}
     \ControlFlowTok{return}
\NormalTok{         expectedReference }\OperatorTok{==}\NormalTok{ current}\OperatorTok{.}\FunctionTok{reference} \OperatorTok{\&\&}
\NormalTok{         expectedStamp }\OperatorTok{==}\NormalTok{ current}\OperatorTok{.}\FunctionTok{stamp} \OperatorTok{\&\&}
         \OperatorTok{((}\NormalTok{newReference }\OperatorTok{==}\NormalTok{ current}\OperatorTok{.}\FunctionTok{reference} \OperatorTok{\&\&}
\NormalTok{           newStamp }\OperatorTok{==}\NormalTok{ current}\OperatorTok{.}\FunctionTok{stamp}\OperatorTok{)} \OperatorTok{||}
          \FunctionTok{casPair}\OperatorTok{(}\NormalTok{current}\OperatorTok{,}\NormalTok{ Pair}\OperatorTok{.}\FunctionTok{of}\OperatorTok{(}\NormalTok{newReference}\OperatorTok{,}\NormalTok{ newStamp}\OperatorTok{)));}
 \OperatorTok{\}}
\end{Highlighting}
\end{Shaded}

References:

\begin{itemize}
\tightlist
\item
  https://docs.huihoo.com/javaone/2006/java\_se/JAVA\%20SE/TS-3412.pdf
\item
  https://fliphtml5.com/tzor/bqxz/basic
\item
  https://www.artima.com/insidejvm/ed2/index.html
\item
  \href{https://wiki.openjdk.java.net/display/HotSpot/Synchronization}{OpenJDK Wiki - Synchronization}
\item
  https://www.zhihu.com/question/53826114
\item
  \href{https://blogs.oracle.com/dave/biased-locking-in-hotspot}{Biased Locking in HotSpot}
\item
  http://gee.cs.oswego.edu/dl/jmm/cookbook.html
\end{itemize}


\end{document}
